%% The following is a directive for TeXShop to indicate the main file
%%!TEX root = diss.tex

%=======================================
\chapter{Specializing to 2 dimensions}
%=======================================
%=======================================
\section{Introduction}
%=======================================
In a lot of ways the 2-dimensional problem looks like the 1-dimensional problem but there are major differences in the inherent complexity of the work that needs to be done. We have the functional equations
\begin{align*}
KQ - K_1Q_1^0-K_2Q_2^0 + K_{1,2} Q_{1,2}^{0,0} &= \ol{a}_1\ol{a}_2\\
K(1)Q_1^0 - K_2(1)Q_{1,2}^{0,0} &= \ol{a}_2 + a_1 t R_1'\\
K(2)Q_2^0 - K_1(2)Q_{1,2}^{0,0}&= \ol{a}_1 + a_2 t R_2'\\
K(1,2) Q_{1,2}^{0,0}&= 1 + a_1 a_2t R_{1,2}'
\end{align*}
Here
\begin{align*}
	R_1 &= S_1^{-1}Q_1^1\\
	R_2 &= S_2^{-2}Q_2^1\\
	R_{1,2} & = \Sf{0,-1}\Qf{0,1}+\Sf{-1,0}\Qf{1,0}+\Sf{-1,-1}\Qf{1,1}
\end{align*}
with
\begin{align*}
	R_1' &= S_1^{-1}Q_1^1 - \ol{x}_2\Sf{-1,-1}\Qf{1,0}\\
	R_2' &= S_2^{-1}Q_2^2 - \ol{x}_1\Sf{-1,-1}\Qf{0,1}\\
	R_{1,2}' & = \Sf{0,-1}\Qf{0,1}+\Sf{-1,0}\Qf{1,0}+\Sf{-1,-1}\Qf{1,1}
\end{align*}
We can write the functional equations in the explicit form
\begin{align}
(1-tS)Q - (c_1-t\ol{x_1}S_1^{-1}) Q_1^0- (c_2-t\ol{x_2}S_2^{-1}) Q_2^0+ (c_1c_2-t\ol{x_1x_2}S_{1,2}^{-1,-1})Q_{1,2}^{0,0}  &= \ol{a_1a_2}\\
(1-ta_1 S_1^0)Q_1^0 - (c_2-ta_1\ol{x_2}S_{1,2}^{0,-1})Q_{1,2}^{0,0} - a_1 t(S_1^{-1}Q_1^1-\ol{x_2}S_{1,2}^{-1,-1}Q_{1,2}^{1,0}) &=\ol{a_2}\\
(1-ta_2 S_2^0)Q_2^0 - (c_1-ta_2\ol{x_1}S_{1,2}^{-1,0})Q_{1,2}^{0,0} -a_2t(S_2^{-1}Q_2^1-\ol{x_1}S_{1,2}^{-1,-1}Q_{1,2}^{0,1}) &=\ol{a_1}\\
(1-a_1 a_2 t \Sf{0,0})Q_{1,2}^{0,0} - ta_1a_2t(\Sf{-1,-1}\Qf{1,1}+\Sf{-1,0}\Qf{1,0}+\Sf{0,-1}\Qf{0,1})&= 1
\end{align}
where the $R_I'$ terms are written on the left to conserve space. We can already see that we are working in a much richer space than in the 1-dimensional case. We'll mostly work on walks with step sets that don't allow staying in one spot ($\Sf{0,0} = 0$). Sometimes we might though.
%=======================================
\section{Investigating relationships between $G,G_1,G_2$ and $G_{1,2}$}
%=======================================

In this work we'll focus on investigating the properties of the function $G(a_1,a_2) = \Qf{0,0}(a_1,a_2)$. Let $G_1 = G(1,a_2)$, $G_2 = G(a_1,1)$ and $G_{1,2} = G(1,1)$. We could also see these as coefficient extractions of $G$ with respect to $b_1$, $b_2$ and $b_1b_2$, i.e. $G_{1,2} = [b_1^0b_2^0] G$. In previous works it has been shown that there is a relationship between these quantities depending on the particular step set $S$. In some cases (like diagonal walks $S = (x_1+\ol{x_1})(x_2+\ol{x_2})$) we have a linear relationship between these quantities. In other cases (like cartesian walks $S = x_1 +\ol{x}_1 + x_2 + \ol{x}_2$) we don't have such a linear relationship but we have a rational one, that is $G$ is a rational function of $G_1$, $G_2$ and $G_{1,2}$. The aim of the followin sections is to investigate when such relationships exist, when se have a relationship of that type and whether we have other cases.

The work can be split into several cases - walks that never leave the origin and ones that leave but never return to the origin. In the first case the relationship is trivially $G-G_1 - G_2 + G_{1,2} = 0$ because $Q = 1$ itself and so $G = 1$. In the second case $Q$ may be more complicated but again $G = 1$ trivially since no walks end at the origin. This case satisfies the same relationship. However, once we move beyond this stage we need to consider the problem more carefully.

In any of the simple cases we can write $BG = 0$ and sometime $BQ = 0$ where $B = B_1 B_2$ from before. We'll investigate the when we'll have a relationship like $B(FQ) = 0$, $B(FG) = 0$ or $B(FG)$ is constant for some $F$'s.

%=======================================
\section{Simple cases: Walks never leaving the origin and walks leaving the origin but never returning}
%=======================================
Here are all the easy cases of step sets that lead to walks that either don't leave the origin or leave it and never come back. We're not allowing resting steps here. There are 80 such sets among the 256 non-resting possible cases. Of course we could use symmetries to eliminate cases but for now this was an easier thing to do.

There are $2^5 = 32$ types of walks that don't leave the origin. Those are walks with steps among the set $\{$NW,W,SW,S,SE$\}$. Next we look at the walks that leave the origin but never come back.  There is the case of $S = \{NE\}$ that never comes back to any axis. Let's look at those walks that never come back to the $x_1$ axis. We have some steps that go in the positive $x_2$ direction, that is one of $\{$N, NE, E$\}$ together with some steps from $\{$NW, W, SW$\}$. If we were to include any other steps we would come back to the $x_1$ axis. If we only have a N step, we have 8 new cases, if we have only a NE step we have 3 new cases and if we have both, we have 4 new cases. This gives a total of 15 new cases. There's an equivalent number of walks that leave the $x_2$ axis and never come back to it. The running total is $32 + 1  + 15 + 15 = 63$.

To finish off we have 17 cases of walks that leave the origin, eventually interact with both axes but never come back to the origin. For this we must only include steps in such a way that we cannot pick a selection of steps that will decrease both the $x_1$ and $x_2$ coordinates. We have to pick a subset of steps from $\{$N, NE, E$\}$ so the only other steps allowed are NW and SE. The new cases coming from $\{$N, NE, E$\}$, $\{$N,E$\}$, $\{$N, NE $\}$, $\{$NE, E$\}$, $\{$N$\}$, $\{$NE$\}$, $\{$E$\}$ are 4, 4, 2, 2, 1, 2, 2 respectively.

Before we consider walks that come back to the origin we'll look at the behavior $Q$ depending on whether $S$ factors in some nice way or not.

%=======================================
\section{Simple case: Factoring $S$ into a product}
%=======================================

There's something we can say about the case where $S$ factors as $S = S_1(x_1)S_2(x_2)$. We have
\[
	 (DS)^k = (D_1D_2S_1S_2)^k = (D_1S_1)^k(D_2S_2)^k
\]
since the $D_1$ and $D_2$ operators commute and the $D_iS_i$ operators also commute since they deal with different variables. The functional equation for $Q$ is
\[
	(1-tDS)Q = 1
\]
We can compute $Q$ directly via the iteration as 
\[
	Q = \sum_{k = 0}^{\infty} t^k (DS)^k = \sum_{k = 0}^{\infty} t^k (D_1S_1)^k(D_2S_2)^k = Q_1(x_1, a_1,t) \odot_t Q_2(x_2,a_2,t)
\]
where $\odot_t$ denotes the Hadamard product of two functions with respect to variable $t$. From now on we'll drop the subscript $t$ unless it is needed. Here $Q_i$ is given by
\[
	Q_i = \sum_{k = 0}^{\infty} t^k (D_iS_i)^k
\]
and by reversing the iteration we obtain the functional equation
\[
	(1-tD_iS_i)Q_i = 1
\]
In other words, $Q_i(x_i,a_i,t)$ is the generating function for a 1-dimensional problem with Laurent step polynomial $S_i(x_i)$. In the case of a factoring of this sort, we are able to glue the 2-dimensional solution out of two 1-dimensional solutions. If we want to look at the generating function for quarter plane walks of length $k$ it is enough to look at the product of two generating functions for one-dimensional walks of length $k$. Also, we can write
\begin{align*}
	Q &= Q_1 \odot Q_2, \quad \wh Q = \wh Q_1 \odot Q_2\\
	\wc Q &= Q_1 \odot \wc Q_2, \quad \wc{\wh Q}  = \wh Q_1 \odot \wc Q_2
\end{align*}
Therefore
\begin{align*}
	B_1 B_2 Q = Q-\wh Q -\wc Q + \wc{\wh Q} &= (Q_1 - \wh Q_1)\odot (Q_2 - \wc Q_2) \\
	B_1B_2(Q_1\odot Q_2) &=( B_1 Q_1) \odot ( B_2 Q_2)
\end{align*}
In general if we are applying operators that don't involve $t$ to a Hadamard product with respect to $t$ we can treat the operator as a constant with respect to $t$ and apply it on whichever term we would like to. If it applies to both terms then we have to apply it to their mixture in a non-trivial way. For example, each of the $Q_i$'s solves a 1-dimensional problem and therefore we have a relationship of the type $B_i (f_i Q_i) = g_i$ for some $f_i$ and $g_i$. Then we can write
\[
	B_1 B_2 (f_1 Q_1 \odot f_2 Q_2) = ( B_1 f_1 Q_1) \odot ( B_2 f_2 Q_2) = g_1 \odot g_2
\]
but so far this is not a relationship of the type $B_1B_2 (fQ) = g$ since taking out the functions $f_1$ and $f_2$ from within the Hadamard product is non-trivial in general. One might be tempted to say that $f = f_1 \odot f_2$ but that would be incorrect. It would be neat to see if there is a way to find $f$ in terms of $f_1$ and $f_2$.

There is a way to see the Hadamard product solution without the iteration method. Suppose that $Q_i$ is a solution to the functional equation $(M_i - tE_iS_i)Q_i = \ol{a}_i$ for $i = 1$ and $i = 2$. Then we claim that $Q = Q_1 \odot Q_2$ is a solution to the functional equation $(M - tES)Q = \ol{a}_1\ol{a}_2$. To this end we calculate
\begin{align*}
	(M-tES)Q &= M (Q_1\odot Q_2) - tES(Q_1\odot Q_2)\\
	& = (M_1Q_1\odot M_2Q_2) - t((E_1S_1Q_1)\odot (E_2S_2Q_2))\\
\end{align*}
Then we substitute $ M_i Q_i = \ol{a}_i + tE_iS_iQ_i$ into the first term to obtain
\begin{align*}
	 (M_1Q_1\odot M_2Q_2) &= (\ol{a}_1+ t E_1 S_1 Q_1)\odot(\ol{a}_2 + t E_2 S_2 Q_2)\\
	 & = \ol{a}_1\odot\ol{a}_2 + \ol{a}_1\odot (t E_2 S_2 Q_2) + \ol{a}_2\odot(t E_1 S_1 Q_1) + (t E_1 S_1 Q_1)\odot(t E_2 S_2 Q_2))
\end{align*}
We have
\[
	\ol{a}_i\odot (t E_j S_j Q_j) = \ol{a}_i [t^0]t E_j S_j Q_j = 0
\]
since we are dealing with power series in $t$. Therefore
\[
	(M_1Q_1\odot M_2Q_2) =  \ol{a}_1\ol{a}_2 +t ((E_1 S_1 Q_1)\odot (E_2 S_2 Q_2))
\]
and upon substitution we are left with $\ol{a}_1\ol{a}_2$ which is what we needed in the first place.

One may also ask whether we can do something in the case where $S$ factors and the terms share variables.

This factoring also generalizes to higher dimensions. If $S$ factors into a product of $k$ factors $S_1,\ldots, S_k$ where no two of which share any variables then the problem of finding $Q$ boils down to solving $k$ lower-dimensional problems for functions $Q_1,\ldots, Q_k$, naming variables appropriately and taking the Hadamard product of the functions. That is
\[
	Q = Q_1 \odot \cdots \odot Q_k
\]
The more lower-dimensional problems that we solve, the more higher dimensional problems we can solve. Although it appears that the interesting higher dimensional problems that can be solved in this manner will have a lower and lower proportion among all higher dimensional problems.
%=======================================
\section{Kernel Group Considerations}
%=======================================
It is important to find more functional equations if we can. One way to do this is by applying a transformation on the pair $(x_1,x_2)$ that leaves as many parts of the functional equation as we can but still adding some information in the form of a new equation. One way to do this is to cosider mapping $x_1$ and $x_2$ to some quantities separately and then consider the full group generated by such maps. Since we usually eliminate the kernel in front of $Q$ it is useful to consider transformations that preserve the kernel.  We'll call this group $G(K)$, although it won't be the full group of rational transformations preserving the kernel. For example $(x_1,x_2) \mapsto (x_2,x_1)$ will not be in the group. Any transformation on $(x_1,x_2)$ that preserves the kernel $K = 1 -tS$ must necessarily preserve the step polynomial $S$. Here
\[
	S = \alpha \ol{x} + \beta + \gamma x
\]
where $x = x_1$ or $x_2$ and $\alpha,\beta,\gamma$ change correspondingly. We consider sending $x \mapsto X$ and solve for $X$ to determine the possible transformations we can perform. That is we have
\begin{align*}
	\alpha \ol{x} + \gamma x &=\alpha \ol{X} + \gamma X\\
	\alpha(\ol{x} - \ol{X}) &= \gamma (X-x)\\
	\alpha\frac{X-x}{Xx} &= \gamma (X-x)
\end{align*}
We have cases
\begin{itemize}
\item $\alpha = 0, \gamma = 0$. The equation is then trivially true and any transformation preserves the kernel although this is not very useful though.
\item $\alpha = 0, \gamma \neq 0$. Then $ X = x$ and only the identity transformation preserves the kernel.
\item $\alpha \neq 0, \gamma = 0$. Here $X = x$ as well.
\item $\alpha \neq 0, \gamma \neq 0$. Here we either have $X = x$ or $X\neq x$ and $X = \ol{x}\alpha/\gamma = \ol{x} S_{-1}/S_1$
\end{itemize}
In conclusion, either every transformation works and we don't have an interesting case or we have the identity transformation plus a possible second one. The last transformation creates a non-trivial group and we can have groups of various orders. Let
\[
	g_1 : x_1 \mapsto \ol{x_1}\frac{S_1^{-1}}{S_1^{1}},\quad g_2 : x_2 \mapsto \ol{x_2}\frac{S_2^{-1}}{S_2^{1}}
\]
whenever the $S$ quantities are non-zero. Note that $g_i(g_i(x_1)) = x_1$, that is they are involutions. We have three cases for each $x$
\begin{itemize}
\item 1: Any transformation can be applied
\item 2: We only have the identity transformation
\item 3: We have two transformations
 \end{itemize}
 We'll combine cases for analysis in the form "case $i.j$" to mean we have case $i$ for $x_1$ and case $j$ for $x_2$.

\begin{itemize}
\item [Case 1.1:] Here we can apply any transformation on both $x_1$ and $x_2$. From the transformation analysis this means that $S_1^{-1} = S_1^1 = S_2^{-1} = S_2^1 = 0$ and thus means that we have no steps that go away from the origin. Therefore $S = \Sf{0,0}$. Thus we have only two models - no steps or a resting step. (I should include little diagrams with arrows like before).
\item[Case 1.2]: Here any transformation applies on $x_1$, so $S_1^{-1} = S_1^1 = 0$ so the model has no steps that step away from the origin in the $x_1$-direction. Since the only transformation on $x_2$ is the identity transformation we either have $S_2^{-1} = 0$ or $S_2^1 = 0$ in which case the model consists of only an up step or a down step, together with a possible resting step. Similarly, the 2.1 case involves a single $x_1$ or $\ol{x_1}$ step with a possible resting step.
\item[Case 1.3:] (and 3.1): This case is similar to the previous one except that now we have two transformations for $x_2$ and this means $S_2^{-1}, S_2^1 \neq 0$. This means we must have either both vertical steps $x_2$ and $\ol{x}_2$ in the 1.3 case or two horizontal steps $x_1$ and $\ol{x}_1$ in the 3.1 case.
\item[Case 2.2:] In this case both transformations are the identity transformations. Therefore exactly one of $S_1^{-1}$ and $S_1^1$ is zero and exactly one of $S_2^{-1}$ and $S_2^1$ is zero. WLOG, assume that $S_1^1 \neq 0, S_2^1 \neq 0$ (the other models can be obtained by a rotation around the origing of 90, 180 and 270 degrees). The possible models are (I should include the drawings): $\{$N,E,NE$\}$, $\{$N,NE$\}$, $\{$E,NE$\}$, $\{$N,E$\}$, $\{$NE$\}$. In other words, the subsets of $\{$N,E,NE$\}$ except the ones going either only in the N direction or only in the E direction. This is a case of group $G(K)$ or order 1.
\item[Case 2.3:] In this case there is only one transformation for $x_1$ and two for $x_2$. This means $S_1^{-1} = 0$ (positive case) or $S_1^{1} = 0$ (negative case). Also $S_2^{-1}, S_2^1 \neq 0$ so there must be at least one step going in the positive $x_2$ direction and one in the negative $x_2$ direction. We have 17 cases (and another 17 including a resting step). Of course in the 3.2 case we have the meaning of $x_1$ and $x_2$ reversed.This is a case of group $G(K)$ of order 2.
\item[Case 3.3:] This is the most interesting case with two transformations for $x_1$ and $x_2$ each. This means that there must be at least one step going in a positive and one going in a negative coordinate for each of $x_1$ and $x_2$. There are a lot of such cases (in the hundreds). Previous results show that the group $G(K)$ has orders 4, 6, 8 and infinity depending on the step set in question.
 \end{itemize}
 
 %=======================================
\section{Solving basic 2D models}
%=======================================

%=======================================
\subsection{Cases 1.1, 1.2, 1.3}
%=======================================
A quick glance at the step models involved shows that we are dealing with essentially one-dimensional problems. This means they can be solved easily explicitly. One can also note that these models are all factorable
\[
	S(x_1,x_2) = 1\cdot S(x_2)
\]
Both $Q_1$ and $Q_2$ are functions we solved for in the one-dimensional case
\[
	Q_1(x_1,a_1,t) = \frac{1}{1-a_1 t},\quad Q_2(x_2,a_2,t) = q(x_2,a_2,t)
\]
for some solution $q$ from that section. Therefore
\[
	Q(x_1,x_2,a_1,a_2,t) = Q_1(x_1,a_1,t) \odot Q_2(x_2,a_2,t) = \frac{1}{1-a_1 t}\odot q(x_2,a_2,t) = q(x_2,a_2,a_1 t)
\]
Similarly in the 2.1, 3.1 cases
\[
	Q(x_1,x_2,a_1,a_2,t) = Q_1(x_1,a_1,t) \odot Q_2(x_2,a_2,t) =  q(x_1,a_1,t) \odot \frac{1}{1-a_2 t}= q(x_1,a_1,a_2 t)
\]
and this completes the solution. What about the relationship with the $B_1B_2$ transformation? If $f(a,t)$ is the function that goes together with $q(x,a,t)$ then the function $F(a_1,a_2,t)$ that goes together with $Q(x_1,x_2,a_1,a_2,t)$ is given by
\[
	F(a_1,a_2,t) = f(a_2,a_1 t)
\]
Then we have $B_2 (FQ) = 0$ and $B_1B_2(FQ) = 0$.

%=======================================
\subsection{Case 2.2 - group of order 1}
%=======================================
A quick glance at the step models reveals that this case is mostly trivial. We have a lot of models that never leave the origin and those have a trivial solution
\[
	Q(x_1,x_2,a_1,a_2) = \frac{1}{1-a_1a_2 \Sf{0,0} t}.
\]
We also have a lot of models that are essentially one-dimensional with a single step going along an axis and some steps going off-axis where the latter steps can never be taken by a random walk in this model. They are essentially cases already covered in the previous section. Suppose that the single step is taken in the positive $x_1$ direction. This is equivalent to the $s_{-1} = 0$ case in 1D with function
\[
	q(x,a,t) = \frac{1-s_0 t}{(1-st)(1-a s_0 t)} = \frac{1-s_0 t}{(1-x_1s_1 t)(1-a s_0 t)}
\]
Then we have
\[
	Q(x_1,x_2,a_1,a_2,t) = Q_1(x_1,a_1,t) \odot Q_2(x_2,a_2,t) =  q(x_1,a_1,t) \odot \frac{1}{1-a_2 t}= q(x_1,a_1,a_2 t)
\]
and therefore
\[
	Q(x_1,x_2,a_1,a_2,t) = \frac{1-a_2 \Sf{0,0} t}{(1-(x_1+\Sf{0,0})a_2 t)(1-a_1a_2 \Sf{0,0}t)}
\]
The interesting models in case 2.2 are those with non-negative steps that we listed. Here
\[
	S = \alpha x_1 + \beta x_2 + \gamma x_1 x_2 + \delta
\]
The functional equation and its zero extractions are
\begin{align*}
(1-tS)Q - c_1 Q_1^0- c_2 Q_2^0+ c_1c_2 Q_{1,2}^{0,0}  &= \ol{a_1a_2}\\
(1-ta_1 (\delta+\beta x_2))Q_1^0 - c_2Q_{1,2}^{0,0} &=\ol{a_2}\\
(1-ta_2 (\delta+\alpha x_1))Q_2^0 - c_1Q_{1,2}^{0,0}  &=\ol{a_1}\\
(1-a_1 a_2 t \delta)Q_{1,2}^{0,0} &= 1
\end{align*}
We solve the problem by a simple back-substitution of the variables.
\begin{align*}
\Qf{0,0} & = \frac{1}{1-a_1a_2\delta t}\\
Q_1^0 & = \frac{\ol{a}_2}{1-t a_2 S_1^0}\left(1+\frac{b_2}{1-a_1a_2 \delta t}\right) = \frac{1 -  a_1 \delta t}{(1-t a_1S_1^0)(1-a_1a_2\delta t)}\\
Q_2^0 & = \frac{\ol{a}_1}{1-t a_1 S_2^0}\left(1+\frac{b_1}{1-a_1a_2 \delta t}\right) = \frac{1 -  a_2 \delta t}{(1-t a_2S_2^0)(1-a_1a_2\delta t)}\\
Q &= \frac{1}{1-tS}(\ol{a}_1\ol{a}_2 + c_1 Q_1^0 + c_2 Q_2^0 - c_1 c_2 Q_{1,2}^{0,0})\\
&=\frac{1}{1-tS}\left (\ol{a}_1\ol{a}_2 + \frac{1}{1-ta_1a_2 \delta}\left(c_1 \frac{1 -  a_1 \delta t}{1-t a_1S_1^0} + c_2 \frac{1 -  a_2 \delta t}{1-t a_2S_1^0} - c_1 c_2 \right) \right)\\
&= \frac{1}{1-tS}\left(\ol{a}_1\ol{a}_2 + \frac{1}{1-a_1a_2\delta t}(c_1 f_1(a_1,t) + c_2 f_2(a_2,t)-c_1c_2)\right)
\end{align*}
where the last line explicitly shows that $Q$ is composed of functions that depend only on one of $a_1$, $a_2$. If we want to find a relation of the type $B(FQ) = 0$ then we must first find a suitable function $F$. First of all the $1-tS$ part doesn't play a role since it doesn't depend on $a_1$ or $a_2$. Then later the variables $c_1$ and $c_2$ will vanish when extracting the coefficients of $b_1^0$ and $b_2^0$. We'll now show that the function
\[
	F(a_1,a_2,t) = 1 - a_1 a_2 \delta t
\]
works exactly. First
\[
B((1 - a_1 a_2 \delta t)\ol{a}_1\ol{a_2}) = B(\ol{a}_1\ol{a}_2 - t \delta) = B(\ol{a}_1\ol{a}_2) - B(t \delta) = \ol{a}_1\ol{a}_2- \ol{a}_2-\ol{a}_1 + 1 = c_1 c_2
\]
Then
\[
	B(c_1 f_1+ c_2 f_2-c_1c_2) = c_1 f_1 + c_2 f_2 -c_1c_2 - c_2 f_2 - c_1 f_1 = -c_2 c_2
\]
Putting the two quantities together along with dividing by $1-tS$ shows that $B(FQ) = 0$. Of course this works for $G = \Qf{0,0}$ too.

%=======================================
\subsection{Case 2.3 - group of order 2}
%=======================================
There are two possibilities here - positive and negative. In the negative case $S_1^1 = 0$ and the steps have a non-positive $x_1$ component. Due to the fact that a lot of steps are not going to be taken by random walks in those models, all of these models reduce to the cases of $S$ consising of a rest step and maybe a single step in a positive direction. This takes care of 34 models (or 68 if we include a rest step). The interesting models are those with a positive $x_1$ component, or in other words, models where $S_1^{-1} = 0$. The calculations for these models become increasingly more complicated. The functional equation and its extracted forms read as
\begin{align*}
(1-tS)Q - c_1 Q_1^0- (c_2 -t\ol{x}_2 S_2^{-1})Q_2^0+ c_1c_2 Q_{1,2}^{0,0}  &= \ol{a_1a_2}\\
(1-ta_1 S_1^0)Q_1^0 - (c_2 - ta_1 \ol{x}_2 \Sf{0,-1})Q_{1,2}^{0,0} &=\ol{a_2}\\
(1-ta_2 S_2^0)Q_2^0 - c_1Q_{1,2}^{0,0}  - a_2 t S_2^{-1} Q_2^1&=\ol{a_1}\\
(1-a_1 a_2 t \Sf{0,0})Q_{1,2}^{0,0} - ta_1 a_2 \Sf{0,-1} \Qf{0,1} &= 1
\end{align*}
This is starting to look somewhat daunting. However, we can identify a subproblem that is essentially one-dimensional and therefore already solved. For this we look at the second and fourth equations with the substitution $a_i = a, c_i = c, x_2 = x, ta_1 = \tau, S_1^0 = T, Q_1^0 (x_2)= q(x)$. These equations become
\begin{align*}
(1-\tau T)q - (c - \tau  \ol{x} T_{-1})q_0 &=\ol{a}\\
(1-a \tau T_0)q_0 - a\tau \Sf{0,-1} q_1 &= 1
\end{align*}
where $q_0 = q(a,\tau)$. This is the one-dimensional problem again. The solution here is $q(x,a,\tau)$ so the original problem has a solution
\[
	Q_1^0(x_2) = q(x_2,a_2,a_1t),\quad Q_{1,2}^{0,0} = q_0(a_2,a_1t).
\]
The reason the method worked in this case is that we have steps that keep a random walk staying on the $x_2$ axis and steps that take it away from the $x_2$ but no steps that have it return to it. This is similar to the case where a random walk leaves the $x_2$ axis never to come back. Essentially the problem splits into a one-dimensional one regarding the $x_2$ axis and a further 2D one once we step off the axis. We re-write the functional equation as
\[
	(1-tS)Q - (c_2 -t\ol{x}_2 S_2^{-1})Q_2^0  = \ol{a}_1\phi(x_2) = \ol{a_1}(\ol{a}_2 + b_1 Q_1^0 -b_1c_2Q_{1,2}^{0,0})
\]
Here we can treat $x_2$ as a constant and solve for a one-dimensional problem with $x_1$ as a variable. We can perform the same analysis as in the the one-dimensional case except where we had $\ol{a}_1$ we would now have $\ol{a}_1 \phi(x_2)$. Therefore we can solve for $Q$ explicitly. However I haven't done this yet so I'll leave this for later. Curiously enough this analysis does not use the fact that we are dealing with a group of order 2 anymore than the one-dimensional case uses the fact that $|G(K)| = 2$ for $K = 1 - t(\ol{x}+x)$. Interestingly we can try re-doing the problem using the group structure. It would surely illustrate the methods we would use later in analysing the higher group size models. I've done it on paper and I don't think I've made any mistakes. However I don't obtain the correct answers. I suspect it's because of the following - $\sigma_1 = 1/(tS_1^1) - S_1^0/S_1^1$ which does not expand as a power series in $t$. Perhaps we just cannot plug this solution into the functional equations. Therefore one would have to resort to using the $\sigma_2$ solution without being able to use the group structure at all. I can certainly solve the problem in the above manner and use the series generated for each part of the solution to test whether the second solution really fails, and if it does test if it fails because of the $\sigma_1$ behavior.

\endinput
